\subsection{The exact exceptional set for the prime-to-residue-characteristic method}

\begin{statement}{A.56}{All exceptional primes give complete independence in the natural joint quotient}{Proved}
Let $\ell\ge5$ be prime, $\Pi=\pi_1(U_\ell/\Q)$, and let $P_q$ be
the maximal prime-to-$q$ quotient of its geometric fundamental group.
Write $\Pi_{P_q}$ for the corresponding arithmetic extension and retain
also the residual Legendre representation modulo $\ell$. If
\[
 q\mid\ell(\ell^2-1),
\]
the resulting natural joint quotient is the entire fibre product
\[
 \Pi_{P_q}\mathbin{\times}_{G_\Q} R_\ell,
 \qquad
 R_\ell=\{(M,g):\det M=\bar\chi_\ell(g)\}
       \subset\mathrm{GL}_2(\F_\ell)\times G_\Q.
 \tag{A56.1}
\]
Consequently a fixed global cusp section of $\Pi_{P_q}$ can be paired
with any continuous global representation
$\rho:G_\Q\to\mathrm{GL}_2(\F_\ell)$ of determinant $\bar\chi_\ell$
to give a global section of this natural quotient. If instead
$q\nmid\ell(\ell^2-1)$, the residual representation already factors
through $\Pi_{P_q}$.

Thus the set $T_\ell$ of A.46 is the exact exceptional set for this
factorization and independence argument. This assertion concerns the
specified quotient; it does not assert that its sections lift to $\Pi$.
\end{statement}
\begin{proof}
Put $S=\mathrm{SL}_2(\F_\ell)$ and
\[
 U(t)=\begin{pmatrix}1&t\\0&1\end{pmatrix},\qquad
 L(t)=\begin{pmatrix}1&0\\t&1\end{pmatrix},\qquad
 W=\begin{pmatrix}0&-1\\1&0\end{pmatrix}.
\]
The upper and lower elementary matrices generate $S$. It is perfect:
with $D=\operatorname{diag}(2,2^{-1})$, commutation with $D$ gives
$U(3t)$ and $L((2^{-2}-1)t)$ respectively. Both coefficients are
nonzero for $\ell\ge5$, so every elementary matrix is a commutator.
The Legendre geometric monodromies $U(2)$ and $L(-2)$ generate $S$.
The Fermat subgroup contains the commutator subgroup of the geometric
tripod group, since its quotient is abelian. Its image therefore still
contains $[S,S]=S$. This verifies the full geometric residual image
on $U_\ell$, as in A.47.

We prove directly that $S$ has no nontrivial finite quotient of order
prime to $q$ for any $q\mid|S|=\ell(\ell^2-1)$. Let $\varphi$ be
a homomorphism to such a finite group.
For $q=\ell$, $\varphi$ kills $U(1)$ and hence every $U(t)$;
conjugation then kills all $L(t)$, so it is trivial.
For $q=2$, it kills $W$, which has order four. The identity
$WU(t)W^{-1}=L(-t)$ then identifies the images of the upper and lower
elementary subgroups. The image is abelian, and perfectness makes it
trivial.

Suppose $q$ is odd and divides $\ell-1$. Choose $a\in\F_\ell^\times$
of order $q$. The matrix $\operatorname{diag}(a,a^{-1})$ is killed;
its commutators with upper and lower elementary matrices, whose
coefficients $a^2-1$ and $a^{-2}-1$ are nonzero, kill all generators.
Finally suppose $q$ is odd and divides $\ell+1$. A norm-one element
of order $q$ in $\F_{\ell^2}^\times$ gives a nonsplit matrix of
determinant one. In a cyclic basis it has the form
\[
 C=\begin{pmatrix}0&-1\\1&t\end{pmatrix}=WU(t).
\]
Changing basis in $\mathrm{GL}_2(\F_\ell)$ is an automorphism of $S$,
so we may use this form. The matrix $C$ is killed. Thus
$\varphi(W)=\varphi(U(-t))$, which commutes with every
$\varphi(U(a))$. The same identity with $W$ gives
$\varphi(L(-a))=\varphi(U(a))$. Again the image is abelian and hence
trivial. These cases exhaust the prime divisors of $|S|$.

The geometric joint image is a closed subdirect product of $P_q$ and
$S$. By the profinite Goursat argument its possible failure to be the
full product is measured by a common quotient. This quotient is finite
because it is a quotient of $S$, and of order prime to $q$ because it
is a quotient of $P_q$. It is therefore trivial. The geometric image
is $P_q\times S$. The determinant formula and surjectivity to $G_\Q$
identify the residual arithmetic quotient with $R_\ell$; over each
$g\in G_\Q$ its matrices form the complete determinant coset of $S$.
The same coset argument using the full geometric product proves
(A56.1).

For a cusp section $\sigma$ of $\Pi_{P_q}$ and any $\rho$ as stated,
the formula
\[
 g\longmapsto\bigl(\sigma(g),(\rho(g),g)\bigr)
\]
is the asserted section. When $q\nmid|S|$, the finite geometric
residual representation has order prime to $q$, so it kills the kernel
of geometric prime-to-$q$ completion and factors through $\Pi_{P_q}$.
Conversely the full geometric product just proved precludes such a
factorization at every exceptional prime. The executed program
\texttt{a\_main\_exceptional\_groups.py} independently enumerates
$S$ for $\ell=5,7,11,13$ and checks that the normal closures of the
indicated $q$-power-order elements are all of $S$. The displayed
argument, rather than these finite checks, proves the uniform result.
\end{proof}

\begin{statement}{A.57}{Explicit failures of every required exceptional-prime condition}{Proved}
For every prime $\ell\ge5$ and every $q\in T_\ell$, there is a
global section of the natural joint quotient in A.56 whose entire
prime-to-$q$ projection is a fixed global cusp section, but for which
no quadratic twist of its residual representation has the following
local property:
\[
 \begin{array}{c|c}
 q&\text{property that fails}\\\hline
 q\ne2,\ell&\text{unramifiedness at }q\\
 q=\ell&\text{finite flatness over }\Z_\ell\\
 q=2&\text{Artin conductor exponent at most four at }2.
 \end{array}
\]
The residual representations used here are actual $\ell$-torsion
representations of explicitly given elliptic curves over $\Q$.
These examples show the insufficiency of the retained quotient data;
they make no claim of a lift to a full global fundamental-group section.
Each named residual condition can also fail with every other condition
in A.58 imposed simultaneously. For this stronger independence claim
we use the explicit global Kummer representations described below.
\end{statement}
\begin{proof}
For odd $q$ put
\[
 E_q:y^2=x(x-1)(x-q),\qquad \rho=\rho_{E_q,\ell}.
\]
Its invariants are
\[
 c_4=16(q^2-q+1),\qquad \Delta=16q^2(q-1)^2.
\]
Thus at $q$ the equation is minimal with multiplicative reduction of
type $I_2$. After an unramified quadratic extension it is a Tate curve
whose parameter has valuation two. If $q\ne\ell$, the residual
inertia image is a nontrivial unipotent group of order $\ell$, since
$2\not\equiv0\pmod\ell$. Restriction to a subgroup of index at most
two still surjects onto this group. The kernel of any quadratic twist
therefore retains nontrivial inertia, proving the first case.

If $q=\ell$, the Tate extension class has nonzero valuation modulo
$\ell$. It is the very ramified Kummer case of
\citet[\S2.4, especially (2.4.7)]{AMSerre1987}. Its Serre weight is
$\ell+1$, so it is not finite flat, by that paper's Propositions 3--4
in \S2.8. An unramified quadratic twist does not change this inertia
calculation. A ramified quadratic character has restriction
$\omega^{(\ell-1)/2}$ to inertia, where
$\omega=\bar\chi_\ell|_{I_\ell}$. The twisted semisimplification is
\[
 \omega^{(\ell-1)/2}\oplus\omega^{(\ell+1)/2}.
\]
For $\ell\ge5$ these are not the characters $1,\omega$ permitted in
the reducible finite-flat case. They also cannot be the fundamental
level-two pair of the irreducible finite-flat case, because both are
$\F_\ell$-valued. Propositions 3--4 therefore exclude every ramified
quadratic twist as well.

For $q=2$ use $E:y^2=x^3-x$. The eight classes
\[
 d\in\{1,-1,2,-2,5,-5,10,-10\}
\]
exhaust $\Q_2^\times/\Q_2^{\times2}$: they represent both valuation
parities and all four odd-unit classes modulo eight. The twist has
equation $E^{(d)}:y^2=x^3-d^2x$, with $c_4=48d^2$ and
$\Delta=64d^6$. The exact Tate algorithm gives
\[
 \begin{array}{c|c|c|c}
 d&v_2(\Delta_{\min})&\text{Kodaira symbol}&f_2(E^{(d)})\\\hline
 1,-1,5,-5&6&\mathrm{III}&5\\
 2,-2,10,-10&12&I_2^*&6.
 \end{array}
 \tag{A57.1}
\]
The supplied \texttt{a\_main\_exceptional\_twists.gp} runs
\texttt{elllocalred} \citep{AMPARI2026} on all eight exact integral equations and records
the full numerical output in its accompanying \texttt{.out} file;
it also checks their invariants and conductor exponents. The integers
are Tate-algorithm results, with no numerical rank computation.

To transfer (A57.1) to residual representations, both tame and wild
invariants must be preserved. Every twist has integral $j=1728$ and
acquires good reduction over a finite extension of ramification degree
dividing twenty-four; see \citet[\S7.1, before Lemma 7.1.3]{AMCDT1999}.
Thus the action on $T_\ell E^{(d)}$ of the inertia group, including
its tame part, factors through a finite group $H$ of order dividing
twenty-four. For $\ell\ge5$ every subgroup $H_i$ of $H$ has invertible
order in $\Z_\ell$. Its averaging idempotent
\[
 |H_i|^{-1}\sum_{h\in H_i}h
\]
shows that taking $H_i$-invariants commutes with reduction modulo
$\ell$. The invariant dimensions for inertia and for every positive
ramification group are unchanged. The tame codimension and the Swan
conductor are therefore unchanged, so
$f_2(E^{(d)}[\ell])=f_2(E^{(d)})\in\{5,6\}$. This proves the
last case for every quadratic twist.

Each chosen representation has determinant $\bar\chi_\ell$, by the
Weil pairing. A.56 pairs it with the same fixed global prime-to-$q$
cusp section. The resulting global quotient section is consequently
an explicit witness to the stated information loss. Full globality of
this quotient section has been built into the construction; the
additional obstruction to lifting through its discarded kernel has
not been solved or assumed away.

For every exceptional prime the other residual hypotheses can all be
imposed at once. At $q=2$, the curve $y^2=x^3-x$ is already
good at every odd prime, so its representation is finite flat at
$\ell$ and unramified outside $2\ell$. At $q=\ell$, use instead the
global Kummer extension
\[
 \rho(g)=\begin{pmatrix}\bar\chi_\ell(g)&c(g)\\0&1\end{pmatrix},
 \qquad [c]=\kappa_\ell(\ell)\in H^1(\Q,\mu_\ell).
\]
It is unramified at every prime other than $\ell$, but its Tate-form
extension class has valuation one at $\ell$, so the same finite-flat
exclusion for all quadratic twists applies. For every odd
$q\ne\ell$, including the primes dividing $\ell+1$, take
\[
 \rho_q(g)=\begin{pmatrix}\bar\chi_\ell(g)&c_q(g)\\0&1\end{pmatrix},
 \qquad[c_q]=\kappa_\ell(q)\in H^1(\Q,\mu_\ell).
\]
It is unramified outside $q\ell$ and has nonzero unipotent inertia
of order $\ell$ at $q$, because $v_q(q)=1$. Every subgroup of index
at most two still maps onto that inertia image, so no quadratic
twist makes it unramified there. It is finite flat at $\ell$ by the
following direct model. Over $\Z_\ell$, form the disjoint union
\[
 \mathcal G_q=\coprod_{i=0}^{\ell-1}
       \operatorname{Spec}\Z_\ell[X]/(X^\ell-q^i)
\]
and use the group law
\[
 (i,x)(j,y)=
 \left(i+j\bmod\ell,\frac{xy}{q^{\lfloor(i+j)/\ell\rfloor}}\right).
\]
All the displayed denominators are units. This is a finite flat
commutative group scheme, in an exact sequence
$0\to\mu_\ell\to\mathcal G_q\to\Z/\ell\to0$.
Choosing $q^{1/\ell}$ in the generic fibre identifies its Galois
representation with $\rho_q$. Thus finite flatness is verified
directly, not inferred merely from the shape of inertia characters.
At two it is unramified. Therefore the named $q$-condition fails
while every other residual condition in A.58 holds. A.56 realizes
these Kummer representations in the same natural global quotient,
with fixed entire prime-to-$q$ cusp projection. The assertion of
actual elliptic origin above concerns the first set of witnesses;
this stronger independence construction needs only the explicitly
given two-dimensional global representations.
\end{proof}

\begin{statement}{A.58}{An explicit sufficient conjunction for the modular Fermat proof}{Proved}
Let $\ell\ge5$ be prime and $s$ a full section of $U_\ell/\Q$.
The following two conditions suffice to make $s$ globally cuspidal.

\textup{(R)} Its residual Legendre representation has a quadratic twist
$\rho_0$ that is unramified outside $2\ell$, finite flat over
$\Z_\ell$, and has Artin conductor \emph{exponent}
$a_2(\rho_0)\le4$ at two.

\textup{(D)} After choosing an invariant line supplied by \textup{(R)},
let $\sigma$ be the induced section on the corresponding connected
component of $Z_\ell$. There is a cofinal family of finite etale
section neighbourhoods $V_i\to Z_\ell$ of $\sigma$ such that each
smooth compactification $C_i$ satisfies
\[
 C_i(\mathbb A_\Q)_\bullet^{\mathrm{f\text{-}ab}}
       \ne\varnothing.
 \tag{A58.1}
\]
No compatible choice of adelic points in (A58.1) is required. Finite
abelian descent is sufficient here; survival of all finite torsors
is a stronger hypothesis. Local point-or-cusp geometricity at the
finite set $T_\ell$ suffices for \textup{(R)}, by A.46.
These are explicit sufficient inputs, not an assertion of logical
minimality among all possible local hypotheses.
The numerical bound four is sharp for the residual reducibility
implication under only these residual hypotheses: replacing it by
five admits absolutely irreducible examples for every $\ell\ge5$.
\end{statement}
\begin{proof}
We first justify the weaker bound in \textup{(R)} independently of the
stronger level-two calculation in A.6. The determinant of $\rho_0$
is $\bar\chi_\ell$, and hence it is odd. If it were irreducible,
complex conjugation with distinct eigenvalues $1,-1$ would make it
absolutely irreducible, by the argument of A.6. Finite flatness gives
Serre weight two \citep[Proposition 4]{AMSerre1987}. The precise Serre
modularity theorem used in A.6 therefore produces a weight-two cusp
form of level dividing $16$; the prime-to-$\ell$ conductor is a power
of two of exponent at most four. Here we use
\citet[Theorem 9.1]{AMKhareWintenbergerI2009} with the missing
hypothesis supplied by \citet[Corollary (0.2)]{AMKisin2009}.
The nebentypus is trivial: its reduction is
$\det(\rho_0)\bar\chi_\ell^{-1}=1$, while a character of
$(\Z/16)^\times$ has power-of-two order, whose reduction in
characteristic $\ell\ge5$ is injective. But the modular curve
$X_0(16)$ has index $24$, six cusps, and no elliptic points, giving
\[
 g(X_0(16))=1+24/12-6/2=0.
\]
Pullback embeds the weight-two cusp forms of every divisor level
into those of level sixteen. The required form cannot exist.
Thus $\rho_0$ is reducible over $\F_\ell$.

Write its semisimplification as
$\alpha\oplus\bar\chi_\ell\alpha^{-1}$.
The finite-flat inertia classification in Serre's Propositions 3--4
implies that $\alpha|_{I_\ell}$ is either $1$ or
$\bar\chi_\ell|_{I_\ell}$: the fundamental level-two case has no
$\F_\ell$-rational invariant line. Both characters are unramified
outside $2\ell$. At two the cyclotomic character is unramified, so
the two character conductors are equal. Moreover
\[
 2a_2(\alpha)
   =a_2(\rho_0^{\mathrm{ss}})
   \le a_2(\rho_0)\le4.
 \tag{A58.2}
\]
For clarity, the inequality does not assume exactness of full inertia
invariants in characteristic $\ell$: the invariant dimension can
only increase under semisimplification. Wild inertia is a pro-two
group, so its finite-image invariants are exact when $\ell\ge5$;
the Swan term is additive. These two facts give (A58.2).

Local reciprocity now makes $\alpha|_{I_2}$ a character of
$\Z_2^\times/(1+4\Z_2)$, which has order two. It is either trivial
or the inertia character of $\Q_2(i)/\Q_2$. Accordingly let
$\varepsilon$ be either the trivial global character or the global
quadratic character of $\Q(i)/\Q$. It is unramified at every odd
prime and kills $\alpha|_{I_2}$. One of
$\varepsilon\alpha$ and
$\varepsilon\alpha\bar\chi_\ell^{-1}$ is now unramified at every
finite prime, including $\ell$. A finite character of $G_\Q$
unramified at every finite prime is trivial by Minkowski's discriminant
theorem. Hence
\[
 (\varepsilon\rho_0)^{\mathrm{ss}}
       \simeq1\oplus\bar\chi_\ell.
 \tag{A58.3}
\]
The extra quadratic twist at two is necessary in this proof with the
bound four; it must not be silently imported from the stronger local
point calculation. Quadratic twists are scalar, so they preserve the
same invariant line and the resulting section $\sigma$ of $Z_\ell$.

Here is the asserted sharpness of that residual cutoff. Take
$E:y^2=x^3-x$. By A.57 its representation $E[\ell]$ has conductor
exponent five at two; it is finite flat at $\ell$ and unramified at
all other primes. If it were reducible, the preceding character
argument would still work, since $2a_2(\alpha)\le5$ again gives
$a_2(\alpha)\le2$. A quadratic twist would have semisimplification
$1\oplus\bar\chi_\ell$. Its full rational two-torsion, together
with the invariant line and, if necessary, the dual $\ell$-isogeny,
would then contradict Mazur's torsion theorem exactly as in A.8.
Hence $E[\ell]$ is irreducible, and its oddness makes it absolutely
irreducible. This is a statement about the residual implication;
no full fundamental-group section with this representation is asserted.

We next isolate precisely what is used from descent. Put
$W=X_0(4\ell)$ and let $\overline Z_\ell\to W$ be the finite map
from A.7 on the chosen component. The curve $W$ has positive genus,
and \citet[Corollary 8.8]{AMStoll2007} gives
\[
 W(\mathbb A_\Q)_\bullet^{\mathrm{f\text{-}ab}}=W(\Q).
\]
In addition, $W(\Q)$ is finite. This follows directly from the
proof of that corollary and its Theorem 8.6: it constructs a nonconstant
map from $W$ to an abelian variety $A/\Q$ with finite $A(\Q)$.
Every fibre of this nonconstant curve map is finite, so the inverse
image of $A(\Q)$ is finite. In particular no unjustified inference
from compactness of a countable space is being made here.

For each $i$, let $f_i:C_i\to W$ be the composite finite map and put
\[
 Z_i=\bigl(f_i^{-1}(W(\Q))\bigr)_{\mathrm{red}}.
\]
This is a finite reduced $\Q$-subscheme of $C_i$.
Choose any point in the nonempty set (A58.1). Functoriality of finite
abelian descent sends it to one diagonal point $w_i\in W(\Q)$.
At every finite place its coordinate therefore belongs to
$Z_i(\Q_p)$. Since $C_i$ maps nontrivially to the positive-genus curve
$W$, it has genus at least one. The density-one formulation of
\citet[Theorem 8.2]{AMStoll2007} gives $Z_i(\Q)\ne\varnothing$.
Only finite places were used for this support condition, so the
real-component correction in \citet[\S3]{AMStollErrata2017} is respected.

The transition maps between section neighbourhoods take $Z_i(\Q)$
into the corresponding earlier finite sets. Any finite collection
has a common refinement in the chosen cofinal family. A point at
that refinement gives a compatible tuple on the collection. Compactness
of the product of the finite sets therefore gives a compatible
rational pro-point. The stabilizer argument in the proof of A.2 now
applies without a Selmer assumption: the geometric tower is cofinal,
the pro-point is fixed by $\sigma(G_\Q)$, and its stabilizer is the
point decomposition group or the cusp decomposition group above its
rational centre. Thus $\sigma$ is point-induced or cuspidal.
The same conclusion holds for its tripod image and then for $s$,
by the finite-cover fixed-coset argument of A.8.

The point alternative is impossible. Its Legendre curve, after the
quadratic twist in (A58.3), has full rational two-torsion and a stable
line of character $1$ or $\bar\chi_\ell$. In the latter case pass to
the dual $\ell$-isogeny, as in A.8. Either case supplies an elliptic
curve over $\Q$ with torsion containing
$\Z/2\oplus\Z/(2\ell)$, contrary to
\citet[Theorem (8)]{AMMazur1977}. Hence $s$ is cuspidal.

For a Selmer section, A.1 supplies (A58.1), and does so with the
stronger all-finite-descent condition. The present argument shows
exactly where that input can be weakened. It does not produce
(A58.1) from geometry at merely finitely many or density-one places.
A.49 proves that such a replacement fails in general for full
sections of a proper genus-one curve; it is not a claimed counterexample
to the unresolved hyperbolic full-section replacement. A.56--A.57
independently delimit what the retained prime-to-$q$ quotient data can
force at each exceptional prime.
\end{proof}
